% CV professionnel (FR) — El Mehdi Haress
% Cible : Analyste Quantitatif Risques de marché (validation de modèles) — Société Générale
% Compilation : xelatex cv-fr-quant.tex
\documentclass[10pt,a4paper]{article}

\usepackage{fontspec}
\usepackage[french]{babel}
\usepackage[T1]{fontenc}
\usepackage[
  top=1.0cm, bottom=1.0cm, left=1.5cm, right=1.5cm
]{geometry}
\usepackage{amsmath}
\usepackage{amssymb}
\usepackage{xcolor}
\usepackage{tikz}
\usepackage{enumitem}
\usepackage{graphicx}
\usepackage{titlesec}
\usepackage{array}
\usepackage[hidelinks]{hyperref}
\usepackage{fontawesome5}

% ---------- Polices ----------
\setmainfont{texgyrepagella}[
  Extension      = .otf,
  UprightFont    = *-regular,
  BoldFont       = *-bold,
  ItalicFont     = *-italic,
  BoldItalicFont = *-bolditalic,
]
\newfontfamily\headingfont{texgyreheros}[
  Extension      = .otf,
  UprightFont    = *-regular,
  BoldFont       = *-bold,
  ItalicFont     = *-italic,
  BoldItalicFont = *-bolditalic,
]

% ---------- Couleurs ----------
\definecolor{accent}{RGB}{28,58,94}
\definecolor{ink}{RGB}{26,26,26}
\definecolor{muted}{RGB}{95,105,120}
\definecolor{rulegray}{RGB}{200,206,214}
\color{ink}

% ---------- Mise en page ----------
\pagestyle{empty}
\setlength{\parindent}{0pt}
\linespread{1.02}
\frenchspacing

\titleformat{\section}
  {\headingfont\bfseries\fontsize{10}{12}\selectfont\color{accent}}
  {}{0em}{\MakeUppercase}[\vspace{-0.9ex}{\color{rulegray}\rule{\linewidth}{0.7pt}}]
\titlespacing*{\section}{0pt}{2.4ex plus 1.4ex minus .2ex}{1.3ex}

% ---------- Gouttière et colonnes ----------
\newlength{\datecol}\setlength{\datecol}{0.165\textwidth}
\newlength{\gutter}\setlength{\gutter}{1em}
\newlength{\bodycol}
\setlength{\bodycol}{\textwidth}
\addtolength{\bodycol}{-\datecol}
\addtolength{\bodycol}{-\gutter}

% ---------- Commandes ----------
% \entry{dates}{intitulé}{organisme — lieu}
\newcommand{\entry}[3]{%
  \par\vspace{1.0ex plus 0.9ex minus 0.1ex}\noindent
  \begin{minipage}[t]{\datecol}
    \vspace{0pt}\raggedright\small\color{muted}#1
  \end{minipage}%
  \hspace{\gutter}%
  \begin{minipage}[t]{\bodycol}
    \vspace{0pt}\raggedright
    {\headingfont\bfseries\small #2}\\[0.15ex]
    {\small\color{accent}#3}
  \end{minipage}%
  \par\vspace{0.3ex}
}

% Puces alignées sur la colonne de droite
\newenvironment{bullets}
  {\begin{list}{\textcolor{accent}{\scriptsize$\blacktriangleright$}}{%
      \setlength{\leftmargin}{\datecol}
      \addtolength{\leftmargin}{\gutter}
      \addtolength{\leftmargin}{1.25em}
      \setlength{\rightmargin}{0pt}
      \setlength{\labelwidth}{0.9em}
      \setlength{\labelsep}{0.35em}
      \setlength{\itemsep}{0.2ex}
      \setlength{\parsep}{0pt}
      \setlength{\topsep}{0.4ex}
      \setlength{\partopsep}{0pt}
      \small}}
  {\end{list}}

% \skillrow{catégorie}{contenu}
\newcommand{\skillrow}[2]{%
  \par\noindent
  \begin{minipage}[t]{\datecol}
    \vspace{0pt}\raggedright{\headingfont\bfseries\footnotesize #1}
  \end{minipage}%
  \hspace{\gutter}%
  \begin{minipage}[t]{\bodycol}
    \vspace{0pt}\raggedright\small #2
  \end{minipage}%
  \par\vspace{0.7ex}
}

% Caractères laissés tels quels : un kern négatif casse l'extraction du mot-clé « C++ » par les ATS.
\newcommand{\cpp}{\mbox{C++}}

\newcommand{\photo}[1]{%
  \begin{tikzpicture}
    \clip[rounded corners=5pt] (0,0) rectangle (2.55,3.264);
    \node[anchor=south west, inner sep=0] at (0,0)
      {\includegraphics[width=2.55cm]{#1}};
  \end{tikzpicture}%
}

\begin{document}

% =====================  EN-TÊTE  =====================
\begin{minipage}[t]{0.76\textwidth}
\vspace{0pt}
{\headingfont\bfseries\fontsize{22}{25}\selectfont EL MEHDI HARESS}\\[0.9ex]
{\headingfont\fontsize{11}{13}\selectfont\color{accent}%
Docteur en mathématiques appliquées \textbar{} Ingénieur CentraleSupélec}\\[0.4ex]
{\small\color{muted}Modélisation stochastique \textbar{} Méthodes numériques probabilistes \textbar{} Mesures de risque}\\[1.4ex]
{\small
\faEnvelope[regular]~\href{mailto:el-mehdi.haress@inria.fr}{el-mehdi.haress@inria.fr}\quad
\faPhone*~+33 7 58 23 41 50\\[0.35ex]
\faLinkedin~\href{https://www.linkedin.com/in/el-mehdi-haress-67750b3aa/}{el-mehdi-haress}\quad
\faGlobe~\href{https://elmehdiharess.org}{elmehdiharess.org}\quad
\faGithub~\href{https://github.com/ElMehdiHaress}{ElMehdiHaress}\\[0.35ex]
\faMapMarker*~Sophia Antipolis \textemdash{} mobilité Île-de-France}
\end{minipage}\hfill
\begin{minipage}[t]{0.20\textwidth}
\vspace{0pt}\raggedleft
\photo{assets/photo.jpg}
\end{minipage}

\vspace{1.2ex}

% =====================  PROFIL  =====================
\section{Profil}
\small
Chercheur quantitatif, \textbf{5 ans d'expérience} en modélisation stochastique : équations différentielles
stochastiques, schémas de discrétisation avec garanties de convergence, calibration sur observations discrètes
et mesures de risque (VaR, CVaR, mesures spectrales). \textbf{Publications} internationales à comité de
lecture, dont plusieurs consacrées à la \textbf{robustesse des modèles} en régime singulier ou à mémoire longue.
Habitué à faire dialoguer preuve mathématique, implémentation et restitution écrite et orale. Objectif :
la \textbf{validation des modèles} et la minimisation du risque.
\normalsize

% =====================  COMPÉTENCES  =====================
\section{Compétences}
\skillrow{Quantitatif}{EDS et EDPS, calcul stochastique et calcul de Malliavin, processus à mémoire
(mouvement brownien fractionnaire), Monte-Carlo, schémas d'Euler et analyse de convergence, inférence
statistique sur processus discrétisés, probabilités en grande dimension.}
\skillrow{Risque \& finance}{Mesures de risque (VaR, CVaR, mesures de risque spectrales), optimisation
sous contrainte de risque, événements rares et queues lourdes, analyse de sensibilité et de stabilité des
modèles, finance stochastique et modélisation du risque (enseigné en M2 à CentraleSupélec).}
\skillrow{Programmation}{\textbf{Python} (NumPy, SciPy, pandas) \textemdash{} usage quotidien depuis 8 ans ;
MATLAB, R, JavaScript/TypeScript, Git, \LaTeX{} ; bases en \cpp{} (formation d'ingénieur).}
\skillrow{Transverse}{Rédaction de rapports techniques et de notes de synthèse, revue critique par les pairs,
présentation devant des publics experts et non experts, encadrement, organisation de conférences
internationales.}
\skillrow{Langues}{Français \textemdash{} langue maternelle. Anglais \textemdash{} courant. Japonais \textemdash{} notions.}

% =====================  EXPÉRIENCE  =====================
\section{Expérience professionnelle}

\entry{09/2026 \textemdash{} auj.}{Ingénieur de recherche \textemdash{} Modélisation stochastique appliquée}{INRIA Sophia Antipolis \textemdash{} équipe CALISTO}
\begin{bullets}
  \item Conception et analyse de méthodes numériques probabilistes pour des systèmes stochastiques
  multi-échelles issus d'applications industrielles.
  \item Implémentation en Python des schémas étudiés et validation numérique des vitesses de convergence
  démontrées théoriquement.
  \item Restitution des résultats sous forme de notes techniques et de présentations auprès d'interlocuteurs
  pluridisciplinaires (mathématiciens, ingénieurs, partenaires industriels).
\end{bullets}

\entry{12/2024 \textemdash{} 08/2026}{Chercheur post-doctorant en analyse stochastique}{University of Leeds (Royaume-Uni) \textemdash{} équipe de K. Dareiotis et K. Lê}
\begin{bullets}
  \item Démonstration de résultats de \textbf{stabilité trajectorielle uniforme} pour des EDS singulières
  dirigées par un mouvement brownien fractionnaire, publiés dans \emph{Stochastic Processes and their
  Applications} (2026).
  \item \textbf{Quantification de la sensibilité} de la dynamique aux perturbations de la spécification et
  des paramètres : cadre directement transposable à l'analyse de robustesse et au risque de modèle.
  \item Co-organisation du workshop international \textbf{FRSSD 2026} (18 intervenants) : sélection et
  coordination des intervenants, conception du site, budget et logistique en lien avec les RH de l'université.
  \item Chargé de travaux dirigés en processus stochastiques (niveau L3, 50 étudiants).
\end{bullets}

\pagebreak

\entry{10/2021 \textemdash{} 11/2024}{Doctorat \textemdash{} Approximation numérique d'EDS et d'EDPS à dérive singulière}{Université Paris-Saclay / CentraleSupélec \textemdash{} dir. L. Goudenège et A. Richard}
\begin{bullets}
  \item Construction de l'un des \textbf{premiers cadres rigoureux d'approximation numérique pour des
  dynamiques stochastiques à dérive distributionnelle} avec garanties de convergence.
  \item Développement d'\textbf{estimateurs conjoints de plusieurs paramètres} (dérive, volatilité, indice de
  Hurst) sur observations discrètes, avec vitesses de convergence prouvées.
  \item Implémentation, validation et publication en \emph{open source} des schémas et estimateurs
  (Python, MATLAB) ; 5 dépôts publics.
  \item \textbf{Enseignement à CentraleSupélec} ($\sim$60 étudiants par an) : intégration et probabilités,
  équations aux dérivées partielles, et \textbf{finance stochastique et modélisation du risque} (M2).
\end{bullets}

\entry{02/2020 \textemdash{} 08/2020}{Stage de recherche \textemdash{} Apprentissage sous contrainte de risque}{Université d'Osaka (Japon) \textemdash{} équipe de M. Holland}
\begin{bullets}
  \item Transposition des \textbf{mesures de risque spectrales} du cadre financier vers l'optimisation
  statistique, avec garanties théoriques sous pertes non bornées et distributions à queues lourdes.
  \item Deux articles acceptés à \emph{AISTATS}, avec implémentation et campagnes numériques associées.
\end{bullets}

\entry{02/2019 \textemdash{} 08/2019}{Stage de recherche \textemdash{} Estimation pour le modèle d'Ornstein--Uhlenbeck fractionnaire}{University of Alberta (Canada) \textemdash{} équipe de Y. Hu}
\begin{bullets}
  \item Estimation \textbf{simultanée de l'ensemble des paramètres} d'un modèle de taux à mémoire longue sous
  observations discrètes.
\end{bullets}

% =====================  FORMATION  =====================
\section{Formation}

\entry{2021 \textemdash{} 2024}{Doctorat en mathématiques appliquées}{Université Paris-Saclay \textemdash{} Gif-sur-Yvette}

\entry{2020 \textemdash{} 2021}{Master 2 «~Mathématiques de l'aléatoire~» \textemdash{} probabilités et statistiques}{Université Paris-Saclay \textemdash{} Orsay}

\entry{2017 \textemdash{} 2021}{Diplôme d'ingénieur \textemdash{} majeure Mathématiques appliquées}{CentraleSupélec \textemdash{} Gif-sur-Yvette}

\entry{2015 \textemdash{} 2017}{Classes préparatoires aux grandes écoles (MP)}{Lycée Pierre de Fermat \textemdash{} Toulouse}

% =====================  PUBLICATIONS  =====================
\section{Publications sélectionnées}
{\footnotesize\color{muted}8 publications au total ; liste complète sur \href{https://elmehdiharess.org}{elmehdiharess.org}.
En probabilités, l'ordre des auteurs est alphabétique.\par}
\vspace{0.6ex}
\begin{itemize}[leftmargin=1.3em, itemsep=0.35ex, topsep=0pt, label={\textcolor{accent}{\scriptsize$\blacktriangleright$}}]
  \footnotesize
  \item K. Dareiotis, \textbf{E.M. Haress}, K. Lê. \emph{Uniform pathwise stability of additive singular SDEs
  driven by fractional Brownian motion}. Stochastic Processes and their Applications, 2026.
  \item L. Goudenège, \textbf{E.M. Haress}, A. Richard. \emph{Numerical approximation of the stochastic heat
  equation with a distributional reaction term}. IMA Journal of Numerical Analysis, 2025.
  \item L. Goudenège, \textbf{E.M. Haress}, A. Richard. \emph{Numerical approximation of SDEs with
  distributional drift}. Stochastic Processes and their Applications, 2024.
  \item \textbf{E.M. Haress}, A. Richard. \emph{Estimation of several parameters in discretely-observed SDEs
  with additive fractional noise}. Statistical Inference for Stochastic Processes, 2024.
  \item M. Holland, \textbf{E.M. Haress}. \emph{Spectral risk-based learning using unbounded losses}.
  International Conference on Artificial Intelligence and Statistics (AISTATS), 2022.
\end{itemize}

% =====================  COMPLÉMENTS  =====================
\section{Informations complémentaires}
\skillrow{Expertise}{Reviewer pour plusieurs journaux scientifiques et Membre de la COST Action CA24104
«~\textsc{Stochastica}~».}
\skillrow{Communication}{Plus de 15 exposés et posters internationaux, dont séminaires invités à
\textbf{Warwick}, \textbf{Rennes (IRMAR)}, \textbf{Évry} et à l'\textbf{Institut Henri-Poincaré}.
Co-organisateur du colloque CJC-MA 2023 (143 participants).}
\skillrow{Encadrement}{Co-encadrement à distance de la thèse de J. Naffrichoux (depuis 2025) sur les méthodes
numériques pour EDPS. Responsable du suivi des tuteurs du dispositif de soutien en mathématiques de
Paris-Saclay (2023).}
\skillrow{Distinctions}{Programme Vivaldi, Fondation Mathématique Jacques Hadamard (2023) ; bourse de
mobilité internationale IDEX, Paris-Saclay (2019).}

\skillrow{Projets personnels}{Conception et développement d'applications web de bout en bout, en autodidacte :
\textbf{Nerduncover}, application de jeu multijoueur, site web pour un salon de beauté \href{https://hanabellaspa.pro}{hanabellaspa.pro}, site web personnel
\href{https://elmehdiharess.org}{elmehdiharess.org}. Compétences acquises : \textbf{React}/TypeScript,
Tailwind CSS, Vite, \textbf{Firebase} (authentification, Firestore), conception de schémas de données.}
\skillrow{Centres d'intérêt}{\textbf{Dessin}, \textbf{photographie}
(\href{https://elmehdiharess.org/photography}{portfolio en ligne}), \textbf{production de musique lo-fi}
(\href{https://www.youtube.com/@sonosketchy}{sonosketchy}), \textbf{volley-ball}.}

\end{document}
